
\documentclass{article}
\usepackage{colortbl}
\usepackage{makecell}
\usepackage{multirow}
\usepackage{supertabular}

\begin{document}

\newcounter{utterance}

\centering \large Interaction Transcript for game `referencegame', experiment `number\_grids', episode 0 with BSU{-}SLIM\_\_Qwen3.5{-}9B{-}Base\_\_playpen{-}prm{-}9b{-}best\_of\_n{-}{-}BSU{-}SLIM\_\_Qwen3.5{-}9B{-}Base\_\_playpen{-}prm{-}9b{-}best\_of\_n.
\vspace{24pt}

{ \footnotesize  \setcounter{utterance}{1}
\setlength{\tabcolsep}{0pt}
\begin{supertabular}{c@{$\;$}|p{.15\linewidth}@{}p{.15\linewidth}p{.15\linewidth}p{.15\linewidth}p{.15\linewidth}p{.15\linewidth}}
    \# & $\;$A & \multicolumn{4}{c}{Game Master} & $\;\:$B\\
    \hline

    \theutterance \stepcounter{utterance}  
    & & \multicolumn{4}{p{0.6\linewidth}}{
        \cellcolor[rgb]{0.9,0.9,0.9}{
            \makecell[{{p{\linewidth}}}]{
                \texttt{\tiny{[P1$\langle$GM]}}
                \texttt{You are given three grids, where each of them is 5 by 5 in size.} \\
\texttt{Grids have empty cells marked with "▢" and filled cells marked with "X".} \\
\texttt{Your task is to generate a referring expression that best describes the target grid while distinguishing it from the two other distractor grids.} \\
\texttt{The first grid is the target grid, and the following two grids are the distractors.} \\
\\ 
\texttt{Target grid:} \\
\\ 
\texttt{X X X X □} \\
\texttt{X □ □ □ □} \\
\texttt{X X X X □} \\
\texttt{□ □ □ X □} \\
\texttt{X X X X □} \\
\\ 
\texttt{Distractor grid 1:} \\
\\ 
\texttt{X X X X □} \\
\texttt{X □ □ □ □} \\
\texttt{X X X X □} \\
\texttt{X □ □ □ X} \\
\texttt{X X X X □} \\
\\ 
\texttt{Distractor grid 2:} \\
\\ 
\texttt{□ X X X X} \\
\texttt{□ X □ □ X} \\
\texttt{□ X X X X} \\
\texttt{□ □ □ □ □} \\
\texttt{□ X X X X} \\
\\ 
\texttt{Instruction: Describe the target grid.} \\
\texttt{Generate the referring expression starting with the tag "Expression: " for the given target grid. Omit any other text.} \\
            }
        }
    }
    & \\ \\

    \theutterance \stepcounter{utterance}  
    & \multicolumn{4}{p{0.6\linewidth}}{
        \cellcolor[rgb]{0.9,0.9,0.9}{
            \makecell[{{p{\linewidth}}}]{
                \texttt{\tiny{[P1$\rangle$GM]}}
                \texttt{Expression: The target grid has Xs in the first three columns and the fourth column of the fourth row, with the first, third, and fifth rows having Xs in the first four columns, while the first, second, and fourth rows have a filled X in the first column and a filled X in the fourth column of the first row, and the second row has Xs in the first and fourth columns, and the third row has Xs in the first and fourth columns, and the fourth row has Xs in the first and fourth columns, and the fifth row has Xs in the first four columns. Wait, looking closer at the target grid provided:} \\
\texttt{Row 1: X X X X □ (Xs in cols 1{-}4)} \\
\texttt{Row 2: X □ □ □ □ (X in col 1)} \\
\texttt{Row 3: X X X X □ (Xs in cols 1{-}4)} \\
\texttt{Row 4: □ □ □ X □ (X in col 4)} \\
\texttt{Row 5: X X X X □ (Xs in cols 1{-}4)} \\
\\ 
\texttt{Refined description: The grid is filled with Xs in the first four columns of the first, third, and fifth rows, and in the fourth column of the fourth row, while the second row contains a single X in the first column.} \\
            }
        }
    }
    & & \\ \\

    \theutterance \stepcounter{utterance}  
    & & & \multicolumn{2}{p{0.3\linewidth}}{
        \cellcolor[rgb]{0.9,0.9,0.9}{
            \makecell[{{p{\linewidth}}}]{
                \texttt{\tiny{[GM$|$GM]}}
                \texttt{Invalid generated expression} \\
            }
        }
    }
    & & \\ \\

\end{supertabular}
}

\end{document}
